% Partie : sets-functions-relations
% Chapitre : functions
% Section : functions-relations

\documentclass[../../../include/open-logic-section]{subfiles}


\begin{document}

\olfileid[fr]{sfr}{fun}{rel}

\olsection{Les fonctions comme relations}

\begin{explain}
Une fonction qui envoie les !!{element}s de~$A$ sur des !!{element}s de~$B$
définit une relation entre $A$ et~$B$ : celle qui relie $x$ à $y$ si et
seulement si $f(x) = y$. On pourrait même, par exemple si l’on souhaite
réduire le nombre de notions fondamentales des mathématiques,
\emph{identifier} la fonction~$f$ à cette relation, c’est-à-dire à un
ensemble de couples. Une question se pose alors : quelles relations
définissent ainsi des fonctions ?
\end{explain}

\begin{defn}[Graphe d’une fonction] Soit $f\colon A \to B$ une fonction.
Le \emph{graphe} de~$f$ est la relation $R_f \subseteq A \times B$
définie par
\[
R_f = \Setabs{\tuple{x,y}}{f(x) = y}.
\]
\end{defn}

\begin{explain}
Par extensionnalité, le graphe d’une fonction est déterminé de manière
unique. De plus, l’extensionnalité des ensembles justifie immédiatement
le principe implicite d’extensionnalité des fonctions : deux fonctions
$f$ et~$g$ de mêmes domaine et ensemble d’arrivée sont identiques si
elles prennent les mêmes valeurs en chaque point de leur domaine.

De même, une relation « fonctionnelle » est le graphe d’une fonction.
\end{explain}

\begin{prop}\ollabel{prop:graph-function}
Soit $R \subseteq A \times B$ telle que :
\begin{enumerate}
\item si $Rxy$ et $Rxz$, alors $y = z$ ; et
\item pour tout $x \in A$, il existe $y \in B$ tel que
$\tuple{x,y} \in R$.
\end{enumerate}
Alors $R$ est le graphe de la fonction $f\colon A \to B$ définie par
$f(x) = y$ si et seulement si $Rxy$.
\end{prop}

\begin{proof}
Supposons qu’il existe $y$ tel que $Rxy$. S’il existait aussi $z \neq y$
tel que $Rxz$, la première condition sur~$R$ serait violée. Ainsi, un tel
$y$, lorsqu’il existe, est unique. La seconde condition en assure
l’existence pour chaque $x \in A$ ; la fonction $f$ est donc bien définie.
On a manifestement $R_f = R$.
\end{proof}

\begin{explain}
Toute fonction $f\colon A \to B$ possède un graphe : une relation
incluse dans $A \times B$, définie par $f(x) = y$. Réciproquement, toute
relation~$R \subseteq A \times B$ qui possède les propriétés de la
\olref{prop:graph-function} est le graphe d’une fonction~$f\colon A \to B$.
Ce lien étroit permet de considérer une fonction simplement comme son
graphe. Autrement dit, on peut identifier les fonctions à certaines
relations, donc à certains ensembles de couples.
\oliflabeldef{sfr:rel:ref:sec}{Il faut toutefois comprendre cette
« identification » dans l’esprit de la \olref[sfr][rel][ref]{sec} :
il ne s’agit pas d’une affirmation sur la métaphysique des fonctions,
mais du constat qu’il est commode de les \emph{traiter} comme certains
ensembles. Cela permet notamment d’envisager}{On peut ainsi envisager}
sur les fonctions des opérations analogues à celles que nous
avons effectuées sur les relations (voir la \olref[sfr][rel][ops]{sec}).
En particulier :
\end{explain}

\begin{defn}\ollabel{defn:funimage}
Soit $f\colon A \to B$ une fonction et soit $C\subseteq A$.

La \emph{restriction} de~$f$ à~$C$ est la
fonction~$\funrestrictionto{f}{C}\colon C \to B$ définie par
$(\funrestrictionto{f}{C})(x) = f(x)$ pour tout $x \in C$.
Autrement dit, $\funrestrictionto{f}{C} =
\Setabs{\tuple{x, y} \in R_f}{x \in C}$.

L’\emph{image} de~$C$ par~$f$ est $\funimage{f}{C} =
\Setabs{f(x)}{x \in C}$. On parle aussi de l’\emph{image directe}
de~$C$ par~$f$.
\end{defn}

\begin{explain}
Ces définitions donnent $\ran{f} = \funimage{f}{\dom{f}}$ pour toute
fonction~$f$.
\oliflabeldef{sfr:rel:ops:sec}{La notion d’image concorde avec celle
de la \olref[sfr][rel][ops]{sec} lorsqu’on identifie les fonctions à
leurs graphes. Pour la restriction, il faut distinguer les deux
conventions.\footnote{Précision éditoriale : l’original annonce un
accord des notions, mais la restriction d’une fonction à $C$ ne
restreint que la première coordonnée, tandis que la restriction d’une
relation définie précédemment intersecte son graphe avec $C^2$.
Ainsi, pour $f(n)=n+1$ sur $\Nat$ et $C=\{0\}$, la restriction
fonctionnelle a pour graphe $\{\tuple{0,1}\}$, alors que la restriction
relationnelle à $C$ est vide. Les deux coïncident lorsque
$\funimage{f}{C}\subseteq C$.}
Deux autres opérations, l’inversion et le produit relatif, demandent
quelques précisions supplémentaires. Nous les donnerons dans la
\olref[inv]{sec} et dans la \olref[cmp]{sec}.}{}
\end{explain}

\end{document}
